On the domain
$\Omega=\{ (\varphi,r) \mid \varphi\in [-\pi,\pi)\wedge r>1\}$
in polar coordinates solve the partial differential equation
\[
\frac{\partial u}{\partial \varphi}
-
r
\frac{\partial u}{\partial r}
=
\sin\varphi
\]
with the boundary conditions
\[
u(\varphi,1) = 0
\qquad\text{for $-\pi \le \varphi < \pi$.}
\]
Note that for polar coordinates, the function $u(\varphi,r)$ needs to
be a $2\pi$-periodic, i.~e.
\[
\lim_{\varphi\to \pi-} u(\varphi,1) = u(\pi,1).
\]

\begin{loesung}
\begin{figure}
\centering
\definecolor{darkred}{rgb}{0.8,0,0}
\definecolor{darkgreen}{rgb}{0,0.6,0}
\begin{tikzpicture}[>=latex,thick,scale=2]
\begin{scope}
\clip (-2,-2) rectangle (2,2);
\fill[color=darkgreen!10] (-2,-2) rectangle (2,2);
\fill[color=white] (0,0) circle[radius=1];
\end{scope}
\begin{scope}
\clip (-2,-2) rectangle (2,2);
\foreach \phinull in {-180,-170,...,180}{
	\draw[color=darkred,line width=1.2pt] plot[domain=0:1.4,samples=20]
		({180*\x/3.1415+\phinull}:{exp(\x)});
}
\end{scope}
\draw[color=darkgreen,line width=1.2pt] (0,0) circle[radius=1];
\draw[->] (0,0) -- (2.3,0) coordinate[label={$r$}];
\fill (0,0) circle[radius=0.025];
\end{tikzpicture}
\end{figure}
This is a quasilinear partial differential equation of first order
which can be solved using the method of characteristics.
The equation for the characteristics with the parameter $t$ are
\begin{equation*}
\left.
\begin{aligned}
\dot{\varphi} &=  1\\
\dot{r}       &= -r\\
\dot{u}       &=  \sin\varphi
\end{aligned}
\right\}
\qquad\Rightarrow\qquad
\left\{
\begin{aligned}
\varphi(t) &= t + \varphi_0 \\
r(t)       &= r_0e^{-t}.     \\
\end{aligned}
\right.
\end{equation*}
The boundary conditions imply that $r_0=1$ and the characteristic
curve starting at the boundary point $(\varphi_0,1)$ has the parametrization
$t \mapsto (\varphi_0+t,e^{-t})$.
These are spirals starting in the unit circle.
Obviously, these curves cover all of $\Omega$ which implies that the
problem is well posed.

What remains is to solve the equation for $u$ and to eliminate the
$\varphi_0$ variable.
The equation is
\[
\dot{u} = \sin(t+\varphi_0)
\]
and can be solved by integrating, giving the solution
\begin{equation}
u(t)
=
\int_0^t \sin(\tau+\varphi_0)\,d\tau
=
\biggl[-\cos(\tau+\varphi_0)\biggr]_0^t
=
-\cos(t+\varphi_0) + \cos\varphi_0.
\label{30000035:u}
\end{equation}
The boundary condition for $t=0$ is
\[
u(0)
=
-\cos(0+\varphi_0) +\cos \varphi_0
=
0
\]
and is therefore automatically satisfied.

We still have to eliminate $\varphi_0$.
Using $t=-\log r$ und $\varphi_0 = \varphi - t = \varphi+\log r$, we get
from
\eqref{30000035:u}
\[
u(\varphi, r)
=
-\cos(-\log r+\varphi+\log r) + \cos(\varphi+\log r)
=
\cos(\varphi+\log r)
-
\cos(\varphi).
\]

We verify that this function really is a solution by computing the derivatives:
\begin{align*}
\frac{\partial u}{\partial \varphi}
&=
-\sin(\varphi+\log r)+\sin(\varphi)
\\
\frac{\partial u}{\partial r}
&=
-
\sin(\varphi+\log r)\frac{1}{r}.
\end{align*}
Substituting them into
\[
\frac{\partial u}{\partial \varphi}
-r\frac{\partial u}{\partial r}
=
-\sin(\varphi+\log r)+\sin(\varphi)
+r
\sin(\varphi+\log r)\frac{1}{r}
=
\sin\varphi,
\]
which is the original equation.
\end{loesung}

\begin{bewertung}
Quasilinear partial differential equation of first order ({\bf Q}) 1 point,
equation for characteristic curves ({\bf C}) 1 point,
solution for $\varphi$ ({\bf \textPhi}) and $r$ ({\bf R}) 1 point each,
solution for $u(t)$, ({\bf U}) 1 point,
elimination of $t$ and $\varphi_0$ ({\bf E}) 1 point.
\end{bewertung}

